\section{Key Findings and Recommendations}

This section distills the most important insights from the multimodal analysis into actionable guidance for researchers working with RSITMD.

\subsection{Key Findings}

\begin{enumerate}
    \item \textbf{Visual-Linguistic Alignment is Strong but Templated.} Captions in RSITMD are tightly coupled with visual content through color words (41.5\% of captions), spatial prepositions (31.3\%), and concrete object nouns. However, all five captions for the same image paraphrase each other within a narrow 10-word template, providing limited linguistic diversity.

    \item \textbf{Severe Class Imbalance Affects Both Modalities.} The training set has a 37.8:1 imbalance ratio (storagetanks: 227 images vs. plane: 6 images). This skews vocabulary statistics toward majority categories and means models without correction will underrepresent minority-class image-text pairs.

    \item \textbf{Compact Vocabulary and Consistent Length Reduce Model Complexity.} With only 3,470 unique words and captions averaging 10.26 words, RSITMD is well-suited for fixed-length text encoding without dynamic padding or adaptive embeddings.

    \item \textbf{Color is the Dominant Visual Feature in Text.} The word \textit{green} alone accounts for 36.8\% of all color word occurrences, reflecting the pervasive role of vegetation in aerial imagery. Color-object relationships provide a reliable modality bridge for cross-modal alignment.

    \item \textbf{Spatial Relationships are Pervasively Annotated.} Spatial prepositions (\textit{near}, \textit{surrounded}, \textit{next to}, \textit{in the middle of}) appear in 31.3\% of captions at the caption level and 52.9\% of images at the image level, offering consistent training signal for spatial reasoning.
\end{enumerate}

\subsection{Recommendations for Multimodal Learning on RSITMD}

\begin{itemize}
    \item \textbf{Address class imbalance.} Use class-weighted sampling or loss functions (37.8:1 ratio), particularly for plane, boat, intersection, and bareland.

    \item \textbf{Use fixed word embedding layers} (4,000--5,000 dimensions) for the 3,470-word vocabulary; no adaptive mechanisms needed.

    \item \textbf{Encode color-object relationships explicitly.} Prioritize color-aware text representations since color adjectives are the primary visual-linguistic bridge.

    \item \textbf{Incorporate spatial reasoning.} Add spatial reasoning modules or contrastive objectives that align spatial phrase embeddings with visual regions.

    \item \textbf{Leverage all five captions per image} in contrastive learning to amplify the training signal and learn semantic equivalence across surface forms.

    \item \textbf{Apply caption augmentation} (e.g., back-translation) to expand linguistic diversity beyond the structured annotation template.
\end{itemize}

\subsection{Summary}

RSITMD provides a solid foundation for cross-modal retrieval in aerial imagery: its compact vocabulary, consistent captioning style, and strong visual-linguistic alignment simplify model design. The main challenges are class imbalance and the templated nature of captions --- both of which can be addressed with appropriate preprocessing and training strategies.
\end{document}
